% ============================================================
%  Chapitre 2 : Étude Fonctionnelle et Technique
% ============================================================

% ---- Page de titre du chapitre ----
\clearpage
\thispagestyle{chaptertitlepage}
\vspace*{\fill}
\begin{center}
\rule{\textwidth}{0.4pt}\\[0.8cm]
{\fontsize{24}{29}\selectfont\bfseries Chapitre 2}\\[0.6cm]
{\fontsize{18}{22}\selectfont Étude Fonctionnelle et Technique}\\[0.8cm]
\rule{\textwidth}{0.4pt}
\end{center}
\vspace*{\fill}
\clearpage

% ---- Contenu du chapitre ----
\chapter{Étude Fonctionnelle et Technique}

\section{Introduction}

Ce chapitre détaille l'étude fonctionnelle et technique de l'application \textbf{Dental Reservation}. Il présente l'architecture globale, les spécifications des besoins fonctionnels et non fonctionnels, ainsi que les choix technologiques pour le backend, le frontend et la base de données.

\section{Architecture Globale}

L'architecture de l'application suit le modèle \textbf{client-serveur à 3 niveaux} tel qu'illustré dans la Figure~\ref{fig:architecture}.

\begin{figure}[H]
\centering
\fbox{\parbox{0.9\textwidth}{\centering\vspace{0.5cm}
\textbf{Frontend (React)} $\xrightarrow{\text{HTTPS / REST}}$ \textbf{Backend (ASP.NET Core 8)} $\xrightarrow{\text{EF Core / SQL}}$ \textbf{SQL Server}\\[0.3cm]
\small Vite + TypeScript + TailwindCSS \hfill JWT Authentication + Swagger \hfill Base CabinetD\\
\small Port 8080 \hfill Ports 7091 / 5057 \hfill Port 1433
\vspace{0.5cm}}}
\caption{Architecture globale de l'application (client-serveur à 3 niveaux)}
\label{fig:architecture}
\end{figure}

\subsection{Flux de données détaillé}

\begin{enumerate}
    \item L'utilisateur interagit avec l'interface React dans le navigateur web.
    \item Le frontend envoie des requêtes HTTP (\texttt{GET}, \texttt{POST}, \texttt{PUT}, \texttt{DELETE}) vers l'API backend. En développement, le serveur Vite agit comme proxy, redirigeant les requêtes \texttt{/api/*} vers \texttt{https://localhost:7091}. En production, Nginx assure ce rôle.
    \item Le middleware d'authentification ASP.NET Core intercepte chaque requête et valide le token JWT dans l'en-tête \texttt{Authorization: Bearer <token>}. Si le token est absent, expiré ou invalide, le serveur retourne \texttt{401 Unauthorized}.
    \item Le contrôleur approprié traite la requête en utilisant \texttt{ApplicationDbContext} (Entity Framework Core) pour construire des requêtes LINQ qui sont traduites en SQL et exécutées sur la base de données SQL Server \texttt{CabinetD}.
    \item La réponse JSON est renvoyée au frontend, qui met à jour l'interface réactive via le système de state management de React et le cache React Query.
\end{enumerate}

\subsection{Justification des choix technologiques}

\begin{itemize}
    \item \textbf{ASP.NET Core 8}~: framework robuste, performant et bien documenté, avec un support natif de JWT et Swagger.
    \item \textbf{Entity Framework Core}~: ORM qui permet d'écrire des requêtes en LINQ au lieu de SQL brut, réduisant les risques d'injection SQL.
    \item \textbf{React 18}~: bibliothèque UI la plus populaire, avec un vaste écosystème de composants et d'outils.
    \item \textbf{TypeScript}~: typage statique pour détecter les erreurs au moment du développement.
    \item \textbf{Vite}~: bundler ultra-rapide avec Hot Module Replacement (HMR) pour un développement fluide.
    \item \textbf{TailwindCSS + shadcn/ui}~: design cohérent et accessible sans écrire de CSS personnalisé.
    \item \textbf{React Query}~: gestion optimisée de l'état serveur avec cache, revalidation et mutations.
\end{itemize}

% ============================================================
\section{Spécifications des besoins}

\subsection{Besoins fonctionnels}

\begin{longtable}{@{} c l p{7.5cm} c @{}}
\toprule
\textbf{\#} & \textbf{Module} & \textbf{Exigence} & \textbf{Priorité} \\
\midrule
\endfirsthead
\toprule
\textbf{\#} & \textbf{Module} & \textbf{Exigence} & \textbf{Priorité} \\
\midrule
\endhead
\bottomrule
\caption{Spécifications des besoins fonctionnels}
\label{tab:besoins_fonctionnels}
\endfoot
BF01 & Auth. & L'utilisateur doit pouvoir se connecter avec un email et un mot de passe & Haute \\
BF02 & Auth. & Le système doit générer un token JWT valide 1 heure après connexion réussie & Haute \\
BF03 & Auth. & Les routes protégées doivent être inaccessibles sans token valide & Haute \\
BF04 & Patients & Le praticien doit pouvoir créer un nouveau patient avec toutes ses informations & Haute \\
BF05 & Patients & Le praticien doit pouvoir modifier les informations d'un patient existant & Haute \\
BF06 & Patients & Le praticien doit pouvoir supprimer un patient & Moyenne \\
BF07 & Patients & Le praticien doit pouvoir rechercher un patient par nom, prénom, téléphone, email ou matricule & Haute \\
BF08 & Patients & Le praticien doit pouvoir consulter le dossier complet d'un patient avec son historique de RDV & Haute \\
BF09 & Patients & Le praticien doit pouvoir imprimer le dossier patient en PDF & Moyenne \\
BF10 & RDV & Le praticien doit pouvoir créer un rendez-vous en sélectionnant un patient, un médecin, un cabinet, une date, une heure, une durée et une nature de soin & Haute \\
BF11 & RDV & Le système doit détecter automatiquement les conflits de créneaux & Haute \\
BF12 & RDV & Le praticien doit pouvoir modifier le statut d'un rendez-vous (confirmé, en attente, annulé) & Haute \\
BF13 & RDV & Le praticien doit pouvoir supprimer un rendez-vous & Moyenne \\
BF14 & RDV & Le praticien doit pouvoir filtrer les rendez-vous par patient, médecin, date ou statut & Moyenne \\
BF15 & Calendrier & Le praticien doit visualiser les RDV sur un calendrier hebdomadaire (08:00--17:30, pas de 30 min) & Haute \\
BF16 & Calendrier & Le calendrier doit afficher un code couleur pour les statuts et un indicateur de l'heure courante & Moyenne \\
BF17 & Créneaux & Le praticien doit pouvoir consulter les créneaux disponibles avec durée configurable & Haute \\
BF18 & Créneaux & Le praticien doit pouvoir réserver directement depuis la grille des créneaux & Moyenne \\
BF19 & Profil & Le praticien doit pouvoir consulter et modifier ses informations de profil & Basse \\
BF20 & Profil & Le praticien doit pouvoir changer son mot de passe & Basse \\
BF21 & Dashboard & Le système doit afficher un tableau de bord avec les métriques clés & Moyenne \\
BF22 & Dashboard & Le tableau de bord doit proposer des accès rapides aux fonctionnalités principales & Basse \\
BF23 & Infos & Le praticien doit pouvoir consulter les informations d'un bénéficiaire/assuré/collectivité & Moyenne \\
BF24 & Attente & Le système doit afficher la liste d'attente d'un médecin pour une date donnée & Moyenne \\
\end{longtable}

\subsection{Besoins non fonctionnels}

\begin{longtable}{@{} c l p{5cm} p{5cm} @{}}
\toprule
\textbf{\#} & \textbf{Catégorie} & \textbf{Exigence} & \textbf{Détails} \\
\midrule
\endfirsthead
\toprule
\textbf{\#} & \textbf{Catégorie} & \textbf{Exigence} & \textbf{Détails} \\
\midrule
\endhead
\bottomrule
\caption{Spécifications des besoins non fonctionnels}
\label{tab:besoins_non_fonctionnels}
\endfoot
BNF01 & Sécurité & Authentification JWT & Validation émetteur, audience, clé HMAC-SHA256, durée de vie \\
BNF02 & Sécurité & Contrôle d'accès & Routes API protégées par \texttt{[Authorize]} \\
BNF03 & Sécurité & CORS & Politique ouverte en développement, configurable en production \\
BNF04 & Performance & Temps de réponse & Requêtes API $<$ 500ms \\
BNF05 & Performance & Cache client & React Query avec invalidation automatique \\
BNF06 & Performance & Compilation rapide & Vite avec HMR $<$ 1 seconde \\
BNF07 & Ergonomie & Responsive design & TailwindCSS (desktop + mobile) \\
BNF08 & Ergonomie & Feedback utilisateur & Notifications toast (succès, erreur, info) \\
BNF09 & Ergonomie & Navigation intuitive & Sidebar rétractable + menu utilisateur \\
BNF10 & Maintenabilité & Séparation & Controllers / Models / DTOs / DbContext \\
BNF11 & Maintenabilité & Typage statique & TypeScript (frontend) + C\# (backend) \\
BNF12 & Déployabilité & Conteneurisation & Docker multi-stage builds \\
BNF13 & Déployabilité & Orchestration & Docker Compose \\
BNF14 & Documentation & API interactive & Swagger UI avec support JWT \\
\end{longtable}

% ============================================================
\section{Spécifications Techniques}

\subsection{Backend --- Stack technique}

\begin{table}[H]
\centering
\caption{Spécifications techniques du backend}
\label{tab:stack_backend}
\begin{tabular}{@{} l l l p{5cm} @{}}
\toprule
\textbf{Composant} & \textbf{Technologie} & \textbf{Version} & \textbf{Rôle} \\
\midrule
Framework & ASP.NET Core & 8.0 (LTS) & Framework web principal \\
ORM & EF Core SqlServer & 9.0.7 & Mapping objet-relationnel \\
Auth. & JwtBearer & 8.0.7 & Middleware JWT \\
Doc. API & Swashbuckle & 6.6.2 & Swagger / OpenAPI \\
BDD & SQL Server & 2019+ & Stockage persistant \\
Runtime & .NET & 8.0 & Environnement d'exécution \\
\bottomrule
\end{tabular}
\end{table}

\subsubsection{Détail des contrôleurs et endpoints}

\begin{longtable}{@{} l l l l p{4.5cm} @{}}
\toprule
\textbf{Contrôleur} & \textbf{Route} & \textbf{Méthode} & \textbf{Endpoint} & \textbf{Description} \\
\midrule
\endfirsthead
\toprule
\textbf{Contrôleur} & \textbf{Route} & \textbf{Méthode} & \textbf{Endpoint} & \textbf{Description} \\
\midrule
\endhead
\bottomrule
\caption{Endpoints de l'API REST}
\label{tab:endpoints}
\endfoot
AuthController & /api/Auth & POST & /login & Authentification JWT \\
\midrule
PatientsController & /api/Patients & GET & / & Liste des patients \\
 & & GET & /\{id\} & Détail d'un patient \\
 & & GET & /stats & Statistiques patients \\
 & & GET & /\{id\}/dossier & Dossier complet \\
 & & POST & / & Création d'un patient \\
 & & PUT & /\{id\} & Mise à jour \\
 & & DELETE & /\{id\} & Suppression \\
\midrule
RendezVousController & /api/RendezVous & GET & / & Liste des RDV \\
 & & GET & /available-slots & Créneaux disponibles \\
 & & POST & / & Création (+ conflits) \\
 & & PUT & /\{num\_rdv\} & Modification \\
 & & DELETE & /\{num\_rdv\} & Suppression \\
\midrule
PraticienController & /api/Praticien & GET & / & Liste des praticiens \\
\midrule
ProfileController & /api/Profile & GET & /me & Profil connecté \\
 & & PUT & /me & Mise à jour profil \\
 & & PUT & /password & Changement mot de passe \\
\midrule
InformationController & /api/Information & GET & /beneficiaire/\{m\} & Infos bénéficiaire \\
 & & GET & /assure/\{m\} & Infos assuré \\
 & & GET & /collectivite/\{code\} & Infos collectivité \\
\midrule
ListeAttenteController & /api/ListeAttente & GET & /by-doctor-date & Liste d'attente \\
\end{longtable}

\subsection{Frontend --- Stack technique}

\begin{table}[H]
\centering
\caption{Spécifications techniques du frontend}
\label{tab:stack_frontend}
\begin{tabular}{@{} l l l p{5cm} @{}}
\toprule
\textbf{Composant} & \textbf{Technologie} & \textbf{Version} & \textbf{Rôle} \\
\midrule
Framework UI & React & 18.x & Composants réactifs \\
Langage & TypeScript & 5.x & Typage statique \\
Bundler & Vite & 5.x & Build + HMR \\
Styles & TailwindCSS & 3.x & Framework CSS utilitaire \\
Composants UI & shadcn/ui (Radix UI) & Dernière & Composants accessibles \\
État serveur & TanStack React Query & 5.56+ & Cache, mutations \\
Routage & React Router DOM & 6.x & Navigation SPA \\
Icônes & Lucide React & 0.462+ & Icônes SVG \\
Dates & date-fns & 3.x & Manipulation de dates \\
Formulaires & React Hook Form + Zod & --- & Validation formulaires \\
\bottomrule
\end{tabular}
\end{table}

\subsubsection{Architecture des pages}

\begin{table}[H]
\centering
\begin{tabular}{@{} l l l p{5.5cm} @{}}
\toprule
\textbf{Page} & \textbf{Fichier} & \textbf{Route} & \textbf{Description} \\
\midrule
Login & Login.tsx & \texttt{/} & Formulaire de connexion \\
Dashboard & Dashboard.tsx & \texttt{/dashboard} & Métriques et accès rapides \\
Calendar & Calendar.tsx & \texttt{/calendar} & Vue hebdomadaire \\
Appointments & Appointments.tsx & \texttt{/appointments} & Gestion CRUD des RDV \\
Patients & Patients.tsx & \texttt{/patients} & Gestion CRUD des patients \\
Patient Dossier & PatientDossier.tsx & \texttt{/patients/:id/dossier} & Dossier + historique \\
Available Slots & AvailableSlots.tsx & \texttt{/slots} & Créneaux + réservation \\
Profile & Profile.tsx & \texttt{/profile} & Gestion du profil \\
\bottomrule
\end{tabular}
\end{table}

\subsection{Structure de la base de données}

\begin{table}[H]
\centering
\caption{Structure de la base de données}
\label{tab:bdd}
\begin{tabular}{@{} l l p{5.5cm} p{3.5cm} @{}}
\toprule
\textbf{Table} & \textbf{Clé primaire} & \textbf{Colonnes principales} & \textbf{Description} \\
\midrule
\texttt{Usercmim} & username & lastName, password, userType & Comptes utilisateurs \\
\texttt{pop\_personne} & Id Num Personne (auto) & NomPer, PrenomPer, DteNaisPer, SexePer, CIN, Email, Phone & Fiches patients \\
\texttt{pop\_adresse\_p} & Id Num Adresse Per & LibAdresse & Adresses \\
\texttt{pop\_contrat\_personne} & Id\_Num\_Personne\_Ctt & --- & Contrats d'assurance \\
\texttt{collectivite} & CodeCollectivite & RaisonSociale, Adresse, Ville & Organismes \\
\texttt{tab\_RDV\_patient} & num\_rdv (auto) & Id\_personne, DateRdv, Heure, NomPs, Status & Rendez-vous patients \\
\texttt{tab\_RDV\_Pra} & Keyless & DateRDV, HeureDebut, NomPra & Planning praticiens \\
\texttt{Praticien\_D} & cod\_pra & NomPra, CodePs & Référentiel praticiens \\
\bottomrule
\end{tabular}
\end{table}

\section{Conclusion}

Ce chapitre a permis de détailler l'architecture globale de l'application, les spécifications des besoins fonctionnels et non fonctionnels, ainsi que les choix technologiques pour chaque couche du système. Le chapitre suivant présentera la conception UML de l'application.
